% =====================================================================
%  Függelék: kiegészítő eredmények (az 5. fejezetből)
% =====================================================================

\chapter{Kiegészítő eredmények}
\label{app-eredmeny-reszletek}

Ez a függelék \az{\ref{ch:eredmenyek}}.~fejezet részletes táblázatait és számait tartalmazza.

% ---------------------------------------------------------------------
\section{A pilot részletei}
\label{app-er-pilot}
% ---------------------------------------------------------------------

A pilot a \code{hu-coffee} szövegen (1\,702 karakter, négy darab) tíz konfigurációval, konfigurációnként három
vizsgával, összesen 300 kérdéssel futott. Ha egy kérdéshely minden kísérlete elbukott, a megtartott kérdésnek
néha csak két-három válasza volt, ami a vizsga $Q$-ját $0{,}25$-tel csökkentette (javítás: a kiválasztás a formát
is figyeli, egy javító hívás pótolja a disztraktorokat); sikertelen hívásoknál egy egykérdéses tartalékvizsga
került a kimenetbe, a kódoló pedig a hosszabb szövegeket a tartalomtól függetlenül vágta
(\ref{sec:csonkolt}.~szakasz). Szivárgásra példa az \arm{E3} egy „kérdése”, amely szó szerint egy forrásmondat
volt („…Európába pedig a 17.~században jutott el a velencei kereskedők révén.”), kulcsa ugyanez volt, disztraktorai pedig ugyanezt a mondatot ismételték a
16., 18. és 15.~századdal; az \arm{E3} kérdéseinek 20\,\%-a szivárgott, 47\,\%-uk nem volt kérdés, a bíró mégis
többnyire 5-ös helyességet adott. Az azonos kontextust kapó \arm{E0} és \arm{E5a} $0{,}83$, illetve $0{,}90$
minőségi indexet ért el, ezért egy dokumentumon a $0{,}07$-nél kisebb különbségek zajnak tekintendők.

% ---------------------------------------------------------------------
\section{Az előzetes vizsgálat helyi konfigurációi}
\label{app-er-elozetes}
% ---------------------------------------------------------------------

\begin{table}[htbp]
\centering
\caption{Az előzetes vizsgálat helyi konfigurációinak eredményei (Qwen2.5-7B, 6 dokumentum).}
\label{tab:elozetes}
\small
\setlength{\tabcolsep}{4pt}
\begin{tabular}{@{}lcccccccc@{}}
\toprule
Konfiguráció & $Q$ & 95\,\% KI & $F$ & $G$ & $A$ & $V$ & Lef. & Min. \\
\midrule
\multicolumn{9}{@{}l}{\emph{Alapvonalak}} \\
\arm{B-doc-L} (teljes dok.)        & 0,83 & 0,76--0,90 & 83           & 79          & 86          & 85          & 45          & 0,71 \\
\arm{E0} (naiv RAG)                & 0,64 & 0,46--0,81 & \textbf{100} & 43          & 62          & 52          & 25          & 0,34 \\
\midrule
\multicolumn{9}{@{}l}{\emph{A javasolt módszer lépései}} \\
\arm{E1} (tervezés)                & 0,88 & 0,82--0,94 & \textbf{100} & 80          & 80          & 91          & 46          & 0,76 \\
\arm{E2} (tervezés + ellenőrzés)   & \cellcolor[HTML]{E6F4EA}\textbf{0,96} & 0,93--0,98 & \textbf{100} & \textbf{90} & \textbf{97} & \textbf{96} & 44          & \cellcolor[HTML]{E6F4EA}\textbf{0,89} \\
\arm{E3} (+ fogalomgráf)           & 0,88 & 0,79--0,95 & \textbf{100} & 83          & 87          & 83          & 40          & 0,68 \\
\midrule
\multicolumn{9}{@{}l}{\emph{Kiegészítő ablációk}} \\
\arm{E2f} (vak próba nélkül)       & 0,90 & 0,84--0,95 & \textbf{100} & 78          & 88          & 92          & 44          & 0,81 \\
\arm{E2h} (hibrid keresés)         & 0,92 & 0,90--0,94 & \textbf{100} & 82          & 92          & 94          & 42          & 0,88 \\
\arm{E2x} (csak szabályok)$^\ast$  & 0,89 & 0,83--0,95 & \textbf{100} & 78          & 83          & 95          & 42          & 0,77 \\
\midrule
\multicolumn{9}{@{}l}{\emph{Darabolási változatok}} \\
\arm{E5a} (régi kódoló)            & 0,87 & 0,81--0,92 & \textbf{100} & 74          & 84          & 90          & 36          & 0,76 \\
\arm{E5b} (magasabb küszöb)        & 0,70 & 0,54--0,86 & \textbf{100} & 39          & 78          & 63          & 35          & 0,45 \\
\arm{E5c} (rögzített darabolás)    & 0,86 & 0,70--0,97 & \textbf{100} & 82          & 82          & 81          & \textbf{51} & 0,50 \\
\bottomrule
\end{tabular}

\vspace{4pt}
\begin{minipage}{\textwidth}
\footnotesize\noindent\textbf{Jelmagyarázat:}
Az egyhívásos konfigurációk három, a tervezettek egy véletlen maggal futottak (a futási időket \az{\ref{tab:koltseg}}.~táblázat, a szerveres konfigurációkat \az{\ref{sec:modellek-szerver}}.~szakasz közli).
$Q$: minőségi index; \textbf{95\,\% KI}: 95\,\%-os bootstrap-konfidenciaintervallum;
$F$, $G$, $A$, $V$: formai megfelelés, forrás-alátámasztottság, vak megoldhatóság és disztraktor-validitás (\%);
\textbf{Lef.}: dokumentum-lefedettség (\%); \textbf{Min.}: a leggyengébb dokumentum $Q$-ja.
A \textbf{félkövér} értékek oszloponként a legjobb eredményt, a \setlength{\fboxsep}{1.5pt}\colorbox[HTML]{E6F4EA}{\textbf{zöld kiemelés}} a legmagasabb átlagos és dokumentumszintű minimum $Q$-t jelöli.
$^\ast$Az \arm{E2x}-ben a második modell ellenőrző hívásai hibával tértek vissza (\ref{sec:fo-komponensek}.~szakasz).
\end{minipage}
\end{table}

Az \arm{E2} a legrobusztusabb (leggyengébb dokumentuma is $0{,}89$), az ellenőrzés az \arm{E1}-hez képest öt
dokumentumon javított, egyen sem rontott ($+0{,}08$), a kiegészítő ablációk mind elmaradtak az \arm{E2}-től, a
naiv RAG pedig az angol szövegeken omlott össze ($0{,}40$; a magyarokon $0{,}88$). Az \arm{E2}--\arm{E0} (hat
javulás) és az \arm{E2}--\arm{E1} (öt javulás, romlás nélkül) korrigálatlanul $p=0{,}031$, Holm-korrekció után
egyik sem szignifikáns (a legkisebb elérhető korrigált érték $0{,}19$).

% ---------------------------------------------------------------------
\section{Az ellenőrző lánc részletei}
\label{app-er-ellenorzo}
% ---------------------------------------------------------------------

\begin{table}[htbp]
\centering
\caption{Az \arm{E2} ellenőrző láncának működése modellenként.}
\label{tab:ellenorzo}
\small
\setlength{\tabcolsep}{4pt}
\begin{tabularx}{\textwidth}{@{}lccccccc>{\raggedright\arraybackslash}X@{}}
\toprule
Modell & $n$ & \makecell{Első\\kísérlet} & \makecell{Kísérlet/\\kérdéshely} & \makecell{Átenge-\\dett} & \makecell{Pontos-\\ság} & \makecell{Felidé-\\zés} & $\kappa$ & Leggyakoribb elutasítási okok \\
\midrule
Qwen2.5-7B     & 60  & 23\,\% & 1,8 & 52\,\% & 0,94 & 0,54 & $0{,}08$ & duplikátum, igaz lehet egy disztraktor, nem kérdés \\
PULI-LlumiX 8B & 85  & 5\,\%  & 2,0 & 19\,\% & 0,56 & 0,21 & $0{,}05$ & duplikátum, forma, nem kérdés, vak eltérés, szivárgás \\
Gemma 3 12B    & 157 & 70\,\% & 1,3 & 95\,\% & 0,88 & 0,95 & $0{,}00$ & duplikátum, vak eltérés \\
\bottomrule
\end{tabularx}

\vspace{4pt}
\begin{minipage}{\textwidth}
\footnotesize\noindent\textbf{Jelmagyarázat:}
$n$: bírói ítélettel rendelkező kérdések száma (\ref{sec:fo-meres}.~szakasz);
\textbf{Első kísérlet}: az első próbálkozásra elfogadott kérdéshelyek aránya;
\textbf{Átengedett}: a végső vizsgában ellenőrzöttként elfogadott kérdések aránya;
\textbf{Pontosság}, \textbf{Felidézés}, $\kappa$: az ellenőrző mint bináris osztályozó mutatói és Cohen-$\kappa$-ja a bíró alátámasztottsági ítéletéhez képest (\ref{sec:statisztika}.~szakasz). Az elutasítási okok az összes kísérletre vonatkoznak.
\end{minipage}
\end{table}

A Qwen ellenőrzője pontos, de szigorú (a jó kérdések közel felét is megjelöli), a Gemmáé szinte mindent átenged;
a PULI mutatói mindössze 85 bírált kérdésen alapulnak. A többi ellenőrzött konfigurációban a $\kappa$ legfeljebb
$0{,}22$ (Qwen \arm{E2h}).

% ---------------------------------------------------------------------
\section{Lefedettség, kognitív szint és felhasználhatóság}
\label{app-er-jelleg}
% ---------------------------------------------------------------------

\begin{sloppypar}
A \emph{lefedettség} az \arm{E0}-ról az \arm{E2}-re a Qwennél 25\,\%$\to$44\,\%, a PULI-nál 10\,\%$\to$21\,\%, a
Gemmánál 24\,\%$\to$43\,\%. A teljes dokumentumos promptnál a lefedettség a Gemmánál 43\,\%, a Qwennél 45\,\%
(tíz kérdésnél a felső korlát $10/K$).

A vizsgaterv a kérdések 43--45\,\%-át a fő vizsgálatban, illetve 27\,\%-át az előzetes vizsgálatban tervezte
\emph{magasabb rendű Bloom-szintre}. A bíró szerint ebből a Gemmánál 7--11\,\%, a Qwennél 5--13\,\%, a PULI-nál
0--2\,\% valósult meg, a tervezett és a megvalósult szint pedig a Gemmánál 12--20\,\%-ban, a PULI-nál
6--10\,\%-ban, a Qwennél 33--45\,\%-ban egyezett.
\end{sloppypar}

A \emph{„felhasználnám”} ítélet aránya (csak a bírált kérdéseken) a Gemmánál a tervezett konfigurációkban
85--88\,\%, a teljes dokumentumos promptnál 83\,\%, az \arm{E5c}-nél 80\,\%, a naiv RAG-nál 42\,\%; a Qwennél
\arm{E2} 80\,\%, \arm{B-doc-L} 68\,\%, \arm{E0} 34\,\%; a PULI-nál legfeljebb 38\,\% (\arm{E5a}, vizsgánként átlagosan
2,3 kérdéssel), az \arm{E2}-ben 29\,\%. A Gemma törzsei a tervezett konfigurációkban átlagosan 122--125
karakteresek (\arm{E3}: 106), a teljes dokumentumos kérdéseknél 63, a naiv RAG-nál 54.

% ---------------------------------------------------------------------
\section{A bírói ítéletek lefedettsége}
\label{app-er-biro}
% ---------------------------------------------------------------------

\begin{table}[htbp]
\centering
\caption[A bírói ítélettel rendelkező kérdések forrásonként.]{A bírói ítélettel rendelkező kérdések
forrásonként (ítélettel rendelkező / összes kérdés) három jellemző konfigurációban.}
\label{tab:biro-lefedettseg}
\small
\begin{tabular}{@{}lcccc@{}}
\toprule
Konfiguráció & Angol (3 dok.) & Saját magyar (3 dok.) & Wikipédia (24 dok.) & Összesen \\
\midrule
Gemma \arm{E2} & 30/30 (100\,\%) & 26/30 (87\,\%)  & 101/240 (42\,\%) & 157/300 (52\,\%) \\
Gemma \arm{E0} & 90/90 (100\,\%) & 90/90 (100\,\%) & 394/720 (55\,\%) & 574/900 (64\,\%) \\
PULI \arm{E2}  & 9/30 (30\,\%)   & 5/30 (17\,\%)   & 71/238 (30\,\%)  & 85/298 (29\,\%) \\
\bottomrule
\end{tabular}
\end{table}

Az áttekintő ítélet a fő vizsgálat tervezett konfigurációinál a kérdések 44--52\,\%-ánál (Gemma), illetve
29--48\,\%-ánál (PULI), az egyhívásosaknál 47--65\,\%-ánál, illetve 70--100\,\%-ánál van meg (az előzetes
vizsgálatban mindenütt 100\,\%); konfigurációnként 85 (PULI \arm{E2}) és 574 (Gemma \arm{E0}) közötti, a tervezetteknél pedig
85--138 (PULI), illetve 132--157 (Gemma) kérdés kapott ítéletet. Hossz szerint nincs érdemi különbség (Gemma
\arm{E2}, Wikipédia- és angol dokumentumok: rövid 46/90, közepes 42/90, hosszú 43/90). Ennek három következménye van:
\begin{enumerate}[label=(\arabic*),leftmargin=*,itemsep=1pt,topsep=2pt]
\item a Wikipédia-dokumentumok $G$, $A$, $V$ értéke kevesebb kérdésen alapul;
\item tartalomfüggő elutasítás esetén e mutatók a bírált részhalmazra érvényesek;
\item ha egy vizsga minden bírói hívása sikertelen, a $Q$ csak a formára (és a vak próbára) támaszkodik, amit az érzékenységvizsgálat ellenőriz (\ref{sec:fo-meres}.~szakasz).
\end{enumerate}
A teljes dokumentumos
konfigurációknál a bíró a dokumentum első 24\,000 karakterét látja (\ref{sec:biro}.~szakasz): a 68\,000 karakteres
jogi szövegen ezért a szervermodell egy helyes kérdését („What is the access economy?”) alá nem támasztottnak
ítélte, mert a fogalom a 38\,603.~karakternél áll.
